\documentclass[11pt]{article}
\usepackage[margin=0.85in]{geometry}
\usepackage{amsmath,booktabs,xurl,hyperref}
\usepackage[T1]{fontenc}
\setlength{\emergencystretch}{2em}
\title{Fourteen-Scenario Test of the Two-Stage Affine Controller}
\author{Pan Dynamics Collision Avoidance Research}
\date{October 1, 2026}
\begin{document}
\maketitle

\section{Result}
All fourteen specified cases were attempted. None completed the requested
avoidance-and-recovery run: all stopped on an unsuccessful optimization stage.
Two cases had strict geometric vehicle overlap before stopping, and twelve
stopped while still separated. An interrupted separated prefix is not a
successful avoidance result. No case reached the nominal-recovery criterion.

Eleven calls stopped in the CLF stage; three stopped in the PCBF stage.
The current version therefore does not pass the scenario experiment. The
previously passing unit checks tested a narrower affine implementation
contract and did not establish these closed-loop outcomes.

\section{Unchanged configuration and measurement}
The frozen controller commit is
\path{905c8bdb9cf4e790e9229133f091d6a05ae13538}. No controller, objective,
constraint or tolerance was changed for this experiment. Seven existing
threat fixtures were run at both 8 and 15 m/s, with exact state observations,
no road-boundary constraints, a 6-mm configured buffer, a 50-ms hold and a
50-ms RK4 prediction step. Primary horizons were 8 and 16 holds, with the
fixed completion allowance producing total lengths of 68 and 76 holds.
The shared soft solve budget remained 5 s.

Every fixture starts with separated rectangles and produces strict overlap
if the ego follows the given path at constant speed. The original known target
motion continues throughout the simulated plant and audit. The plant uses
independent ODE45 integration with relative tolerance $10^{-11}$ and absolute
tolerance $10^{-12}$. An independent Python rectangle implementation checks
31 samples per hold; an additional separating-axis overlap calculation
confirms the two collision cases are strict overlap, not mere touching.
This is sampled geometry, not a continuous-time physical-vehicle certificate.

Each case allowed 1600 holds (80 s), stopping early after five seconds within
all five nominal-recovery tolerances, beginning no earlier than 8 s. In fact,
all cases failed before 8 s, so extending the requested observation cap would
not continue them without addressing the optimizer failure. The recovery
thresholds are 0.1 m lateral error, one degree heading error, 0.1 m/s speed
error, 0.05 m/s lateral velocity error and 0.01 rad/s yaw-rate error.

Two discarded no-target controller calls at each speed warm MATLAB and the
optimizer. Their durations are recorded separately in \path{warmup.json}.
Campaign timings include the entire controller call, including initialization,
assembly and both solves. They exclude plant integration, offline geometry,
file export, MATLAB startup and the separate diagnostic calls. No profiler
was enabled. MATLAB R2026a Update 3 ran with \texttt{-singleCompThread}.

\section{Per-case results}
\begin{center}\small
\begin{tabular}{rlrrrrl}\toprule
Speed & Scenario & Holds & Stop (s) & Min gap (m) & Max call (ms) & Stop stage \\\midrule
8 & headOn & 29 & 1.45 & 2.06438 & 185.7 & CLF \\
8 & acceleratingHeadOn & 8 & 0.40 & 12.79200 & 107.2 & CLF \\
8 & brakingLead & 2 & 0.10 & 1.44229 & 345.3 & CLF \\
8 & crossing & 28 & 1.40 & 2.32047 & 125.9 & PCBF \\
8 & turningCrossing & 6 & 0.30 & 12.45971 & 285.7 & CLF \\
8 & curvedHeadOn & 27 & 1.35 & 2.85544 & 121.5 & PCBF \\
8 & curvedCrossing & 5 & 0.25 & 13.13234 & 402.2 & CLF \\\midrule
15 & headOn & 33 & 1.65 & 0 (overlap) & 129.3 & CLF \\
15 & acceleratingHeadOn & 43 & 2.15 & 0 (overlap) & 165.3 & CLF \\
15 & brakingLead & 4 & 0.20 & 1.41790 & 214.2 & CLF \\
15 & crossing & 152 & 7.60 & 0.47749 & 192.4 & CLF \\
15 & turningCrossing & 59 & 2.95 & 0.04705 & 185.7 & CLF \\
15 & curvedHeadOn & 2 & 0.10 & 29.50636 & 64.3 & PCBF \\
15 & curvedCrossing & 2 & 0.10 & 23.17272 & 183.1 & CLF \\\bottomrule
\end{tabular}
\end{center}
Max call includes the failed controller invocation as well as returned holds.
No plant hold is executed after the failure. The two overlapping simulations
continue kinematically until optimization failure; impact forces are not
modeled. Their prefixes are already failures at first overlap.

\section{Runtime}
There are 400 returned holds and 14 failed calls. Every returned hold invokes
exactly two optimizations. The returned-hold median is 90.107 ms and maximum
is 165.287 ms. Of those 400 holds, 114 exceed 100 ms; excluding each scenario's
first hold still leaves 107 such overruns. Ten failed calls also exceed 100 ms.
Thus 124 of 414 attempted campaign controller calls exceed 100 ms.

The maximum across successful and failed calls is 402.225 ms, at the
8-m/s curvedCrossing failure at simulation time 0.25 s. This is a warmed
runtime call, not MATLAB startup. No campaign call consumes the five-second
budget, but the requested 100-ms maximum is still unmet.

The longest returned call is 15-m/s acceleratingHeadOn frame 21 at 1.00 s:
50.122 ms PCBF solve, 106.644 ms CLF solve, and 8.521 ms remaining controller
work, totaling 165.287 ms. These are the stage timers recorded in the actual
campaign, without a profiling rerun.

\section{Failure diagnostics and collision evidence}
The fourteen saved failure states were replayed once using cache-only renamed
copies that expose the failed search result. All fourteen reproduced their
original stage failures. The eleven CLF cases first returned PCBF exit flag 1,
then CLF flag $-7$; the three PCBF cases returned flag $-2$. No additional
optimization, acceptance logic or controller change was introduced. Diagnostic
timings are separate from campaign runtime results.

Two representative failures were then replayed with the optimizer output
structure captured. In the 8-m/s crossing case, the primary solver reports
``Problem is infeasible.'' This is a verdict on the current affine subproblem,
not proof of physically unavoidable collision.

In 8-m/s curvedCrossing, the primary solve completes in 9 iterations with
PCBF optimum $6.472333\times10^{-9}$. The secondary cap is
$1.006472333\times10^{-6}$. After 77 iterations the CLF solver reports a small
search direction without satisfying its constraint and/or optimality tolerance.
Its primal feasibility measure is $4.888257\times10^{-14}$, dual feasibility
measure 0.146013 and duality gap $6.958170\times10^{-17}$. The evidence in this
representative case identifies failed optimality convergence despite small
primal residual; it must not be described as lack of a feasible trajectory.
The causal source of that numerical conditioning, and which hard rows cause
the three primary infeasibilities, have not been isolated by this experiment.

The first strict overlap samples occur at 1.098333 s in 15-m/s headOn and
1.050000 s in 15-m/s acceleratingHeadOn. The controls covering those samples
have PCBF slack sums 0.0234122 and 0.00658919, respectively, well above the
$10^{-5}$ zero-slack reporting tolerance. Their reported raw affine hard
residuals are $2.695\times10^{-7}$ and $4.729\times10^{-8}$. Therefore these
are executed relaxed affine solutions, not zero-slack safety witnesses.
Observed maximum overlap depths are 0.079246 and 0.069688 m. The measurements
do not isolate the separate contributions of positive slack, linearization
error and prediction/plant integration mismatch.

\section{Artifacts and limits}
\path{report/TWO_STAGE_SCENARIOS_20261001/} contains all fourteen summaries,
all 400 compact frame records, independent audit results, failure-stage and
solver-output diagnostics, actual logs, warm-up records and a reproduction
helper. The manifest records the unchanged source/native hashes and hashes of
the original JSON/MAT exports retained outside the repository. Generated
native binaries are excluded from the commit. No algorithm fixes or claims of
successful recovery accompany this experiment.
\end{document}
